\section{Численно вычисляемые условия границ M3}

Пусть $X=(r_1,\omega_1,r_2,\omega_2,z_1,z_2)$ и $p=(A,B,\nu,\beta)$.
Полный вектор стационарных средних балансов обозначим
\[
F(X;p)=0,\qquad F:\mathbb{R}^6\times\mathbb{R}^4\to\mathbb{R}^6,
\]
где каждый вызов $F$ включает сошедшееся решение внутреннего Fourier
harmonic-balance уравнения $G(\xi;X,p)=0$.

\subsection{Нулевая средняя частота}

Граница обращения направления $i$ находится из расширенной системы
\[
\begin{bmatrix}F(X;p)\\ \omega_i\end{bmatrix}=0.
\]
Это алгебраическое событие не предполагает сингулярность $F_X$ и не
является само по себе бифуркацией устойчивости.

\subsection{Стационарная складка}

Кандидат складки удовлетворяет
\[
F(X;p)=0,\quad F_X(X;p)v=0,\quad \|v\|_2=1,
\]
с дополнительным условием выбора управляющего параметра. В вычислениях
следует контролировать минимальное сингулярное число масштабированного
$F_X$ и подтверждать изменение ветви методом псевдодлины дуги. Отказ
нелинейного решателя не заменяет эти условия.

\subsection{Локальная потеря устойчивости редуцированного потока}

Если согласованный slow-flow имеет вид $\dot X_s=\mathcal{F}(X_s;p)$,
локальная граница определяется
\[
\mathcal{F}(X_s;p)=0,\qquad
\max_j\Re\lambda_j\!\left(D_{X_s}\mathcal{F}\right)=0.
\]
Нулевое вещественное собственное значение и пересечение мнимой оси
комплексной парой требуют разных бифуркационных тестов. Полученная
граница относится к slow-flow, а не автоматически к полной
двухчастотной RF-системе.

\subsection{Геометрические и модельные ограничения}

Геометрические ограничения задаются отдельно, например
\[
r_i+a_i/R < r_{\rm aperture}/R,\qquad
\|\mathbf r_1-\mathbf r_2\|>(a_1+a_2)/R.
\]
Потерю применимости усреднения следует диагностировать через спектр
внутреннего Fourier-оператора, амплитуду соответствующей моды,
unresolved slow/fast residual, cutoff sensitivity и отношение времён.
Эти условия не сводятся к единой алгебраической бифуркационной поверхности.
